%-----------------------------------------------------------
%   Capítulo 4 - Arquitectura do Sistema
%-----------------------------------------------------------
\chapter{Arquitectura de Alto Nível}
\label{cap:arquitectura}

A arquitetura preliminar representa um \textit{gateway} embebido
multiprotocolo com interfaces I2C, RS232, RS485, CAN e Ethernet. As funções
de leitura e escrita articulam-se com o processamento da informação
recebida e com o encaminhamento para outra interface, segundo operações
e mapeamentos definidos. A Figura~\ref{fig:arquitetura-preliminar}
apresenta o percurso funcional de uma operação de encaminhamento.

\begin{figure}[htbp]
    \centering
    \small
    \fbox{\parbox{0.80\linewidth}{\centering
        \textbf{Interface de origem}\\
        I2C, RS232, RS485, CAN ou Ethernet}}
    \par\smallskip
    $\downarrow$
    \par\smallskip
    \fbox{\parbox{0.80\linewidth}{\centering
        \textbf{Módulos de interface}\\
        Leitura / receção}}
    \par\smallskip
    $\downarrow$
    \par\smallskip
    \fbox{\parbox{0.80\linewidth}{\centering
        \textbf{Processamento / normalização}\\
        Segundo a operação definida}}
    \par\smallskip
    $\downarrow$
    \par\smallskip
    \fbox{\parbox{0.80\linewidth}{\centering
        \textbf{Decisão de encaminhamento}\\
        Condicionada ao mapeamento e ao destino}}
    \par\smallskip
    $\downarrow$
    \par\smallskip
    \fbox{\parbox{0.80\linewidth}{\centering
        \textbf{Módulos de interface}\\
        Escrita / transmissão}}
    \par\smallskip
    $\downarrow$
    \par\smallskip
    \fbox{\parbox{0.80\linewidth}{\centering
        \textbf{Interface de destino selecionada}\\
        Conforme a operação de encaminhamento}}
    \caption{Percurso funcional de encaminhamento no gateway multiprotocolo.}
    \label{fig:arquitetura-preliminar}
\end{figure}

Os módulos de interface asseguram a ligação aos dispositivos ou redes
externas e as operações de receção e transmissão de cada tecnologia.
A sua representação em dois pontos do diagrama distingue as funções
de entrada e saída; não impõe a duplicação de componentes nem uma
organização física concreta.

O bloco de processamento interpreta e transforma a informação segundo
a operação definida. Quando necessário, pode efetuar normalização,
preservando os dados e metadados exigidos pela operação de destino.
Esta função não impõe uma representação interna comum a todas as
tecnologias nem uma estrutura de dados específica.

A decisão de encaminhamento determina se a informação processada pode
ser enviada para o destino previsto, de acordo com o mapeamento
estabelecido e com as características das interfaces e dos protocolos
envolvidos. Nem todas as mensagens podem necessariamente ser
encaminhadas para todas as interfaces. Na ausência de um mapeamento
válido, a informação não é encaminhada para esse destino. As operações
de leitura e processamento também podem existir sem transmissão
posterior.

O percurso representado não pressupõe uma correspondência de uma trama
de entrada para uma trama de saída. Conforme a operação e os protocolos
selecionados, podem ser necessários delimitação de tramas
(\textit{framing}), segmentação, armazenamento temporário
(\textit{buffering}), controlo de fluxo e interpretação de protocolos
superiores. A distribuição concreta destes mecanismos pelos blocos
será definida na fase de projeto.

Esta representação é funcional e mantém em aberto a plataforma de
execução, os componentes, a organização interna dos dados e os recursos
de armazenamento. Os pares de interfaces ou protocolos e os mapeamentos
a demonstrar serão concretizados após validação dos requisitos.
